\documentclass[12pt,a4paper]{article}
\usepackage[utf8]{inputenc}
\usepackage[T1]{fontenc}
\usepackage{amsmath}
\usepackage[left=2.00cm, right=2.00cm, top=2.00cm, bottom=2.00cm]{geometry}
\usepackage{amssymb}
\usepackage[italian]{babel}

\title{Soluzione Tutorato 5\\Complementi di Matematica per le Scienze Chimiche\\ A.A. 2021/2022}
\author{prof. Emanuele Bottazzi}
\date{}

\usepackage{amsthm}
\theoremstyle{definition}
\newtheorem{exe}{Esercizio}
\newtheorem*{sol}{Soluzione}

\begin{document}
\maketitle

\begin{exe}
    Calcola integrali multipli.
\end{exe}
\begin{sol}
    \begin{enumerate}
        \item $I_1 = \int_1^2 \frac{1}{x} dx \int_{-2}^{-1} y^2 dy = [\ln x]_1^2 [\frac{y^3}{3}]_{-2}^{-1} = \ln(2) \frac{7}{3}$.
        \item $I_2 = e - 3/2$.
        \item $I_3 = 2^4 = 16$.
    \end{enumerate}
\end{sol}

\begin{exe}
    Definisci insieme $y$-semplice...
\end{exe}
\begin{sol}
    $E = \{(x,y) \in \mathbb{R}^2 : y \in [0,1], y \leq x \leq -y+2 \}$.
\end{sol}

\begin{exe}
    Calcola integrali doppi.
\end{exe}
\begin{sol}
    (1) $1/12$; (2) $-2/3$.
\end{sol}

\begin{exe}
    Calcola $\iint_{A \cup B} 5f(x,y) \, dx dy$.
\end{exe}
\begin{sol}
    $\iint_{A \cup B} 5f = 5 \left( \iint_A f + \iint_B f \right) = 5(1 - 1) = 0$.
\end{sol}

\begin{center}
    prof. Emanuele Bottazzi\\
    emanuele.bottazzi@unipv.it
\end{center}

\end{document}
